%% sections/discussion.tex — Step B10
\section{Discussion}
\label{sec:discussion}

In this study, we extended the DyAM multimodal attention framework with two
complementary contributions: a competitive (OvO) attention mechanism and an
NLP-based encoding of tabular clinical features, and evaluated both using
binary classification and survival endpoints.

The best-performing model, DyAM Rad+IHC-G+Gen+PDL1+NLP~raw, achieved
AUC~$= 0.813$ [95\%~CI: 0.753--0.874], placing it at the
``good-to-excellent'' discrimination threshold ($\text{AUC} > 0.80$) commonly
used in oncology decision support \cite{Vanguri2022}.
For context, this exceeds the single-source models evaluated in the same
cohort under the DyAM framework --- PD-L1~TPS alone (AUC~$= 0.720$) and TMB
alone (AUC~$= 0.616$) --- consistent with the modest discrimination reported
for these markers in the literature \cite{Rizvi2015,Reck2016}.
This confirms that multimodal fusion, augmented with NLP-encoded clinical
context, approaches the upper range of discrimination achievable with
currently available biomarkers in this disease setting.

Why does NLP encoding outperform raw numeric clinical features?
We propose that sentence embeddings capture non-linear interactions between
clinical variables --- for example, the combined prognostic implication of
advanced age, low albumin, and poor performance status --- that are not
explicitly represented in a 13-dimensional numeric vector.
The MiniLM-L6-v2 embedding space \cite{Reimers2019,Wang2020MiniLM} is
pretrained to place semantically similar sentences close together; patients
with similar overall clinical profiles are therefore mapped to nearby
embedding vectors even when no single numeric feature dominates.
Raw numeric encoding, by contrast, presents each variable to the attention
mechanism independently, requiring the model to learn these interactions
de novo from a relatively small training set.

The AUC improvements associated with NLP encoding ($+0.017$ to $+0.030$)
did not reach statistical significance under DeLong confidence interval
testing, as all CIs overlapped substantially.
This is best interpreted as a power limitation rather than an absent effect.
With $n = 247$, the 95\%~CI width for AUC is approximately $\pm 0.07$;
detecting a true $\Delta\text{AUC} \approx 0.02$ at 80\% power would require
approximately 950--1{,}250 patients (Supplementary Note~S3).
The directional consistency of the improvement across both the IHC-A and
IHC-G arms, and across raw and PCA-compressed NLP variants, supports a true
but modest effect size that is currently underpowered rather than null.

The survival analysis provides an important reframing of this finding.
While binary AUC improvements were modest and non-significant, multivariate
Cox regression showed that NLP-PCA16 achieved the highest hazard ratio in the
IHC-A arm (HR~$= 6.07$ vs.\ HR~$= 5.30$ for the no-clinical baseline), and
NLP raw achieved the highest hazard ratio in the IHC-G arm (HR~$= 4.97$ vs.\
HR~$= 4.74$), whereas numeric clinical labs provided no comparable
improvement in either arm.
Furthermore, the time-dependent AUC advantage of NLP encoding was
concentrated at 12 and 18~months ($+0.024$ and $+0.020$ in the IHC-A arm)
rather than at 6~months.
Taken together, these results suggest that NLP encoding of tabular clinical
features predominantly captures \emph{prognostic (time-to-event)} information
rather than \emph{predictive (binary response)} information --- a
distinction with direct implications for how such encodings should be
evaluated in future multimodal studies.

The OvO competitive attention mechanism matched the accuracy of the original
cooperative attention while using a substantially simpler parameterization.
An initial single-partition screen suggested a context-dependent pattern,
with OvO favoured in low-modality configurations
(up to $+3.27\%$ for PDL1+Gen) and the original model favoured as more
modalities were added; the overall best-performing single-partition
configuration (Rad+IHC-G+Gen+PDL1, AUC~$=0.8003$) used OvO attention.
A confirmatory repeated-seed analysis (5~seeds~$\times$~10-fold) placed the
two attention mechanisms within 0.011~AUC of each other at all four primary
benchmarks ($p \geq 0.17$), with OvO reaching a repeated-seed mean AUC of
$0.773$ versus $0.760$ for the cooperative model at the primary
configuration --- directionally favouring OvO, though the single-partition
``best'' configuration's rank relative to the cooperative model reversed
once averaged across seeds, underscoring the value of repeated-seed
evaluation over a single split.
Across all 21 configurations the repeated-seed data also reversed the
\emph{direction} of the context-dependence seen in the screen: OvO matched
or exceeded cooperative attention in 7 of 8 configurations fusing five or
more sources, but in only 5 of 9 configurations with two to four sources,
where the spread between the two mechanisms was several times larger.
A competitive one-vs-others weighting is plausibly better conditioned when
many partially redundant sources compete, whereas with few sources the
comparison against the mean of the remaining modalities is estimated from
one or two values and becomes unstable; we note, however, that the effect
sizes involved ($\leq 0.018$~AUC) sit below the seed-to-seed variability, so
this remains a hypothesis generated by the data rather than a tested claim.
When combined with NLP-clinical encoding, OvO reproduced the identical
fusion-dependent improvement pattern seen in the primary NLP-clinical
analysis (improvement in the same 15 of 21 configurations, decline in the
same six single-modality configurations), reaching AUC~$=0.782$ at the
primary benchmark under the repeated-seed protocol (repeated-seed
NLP-clinical confirmation: $0.783$) and
AUC~$=0.8133$ under the single-partition protocol at the
Rad+IHC-A+Gen+TMB+PDL1 configuration --- the highest value observed under
either protocol in this study.
OvO's practical value lies in its simpler parameterization ---
$N$ independent per-modality scores compared against the mean of the
remaining modalities, rather than the $N \times N$ weight matrix required by
cooperative attention --- delivering matching accuracy, and its improvement
pattern under NLP-clinical encoding generalises across both attention
mechanisms, indicating that the clinical-encoding intervention, not the
choice of attention mechanism, is the source of the observed gain.

From a clinical implementation perspective, NLP encoding of structured
clinical data offers practical advantages: encoding a patient profile
requires less than one second on CPU, requires no labelled clinical text,
and can be applied to any structured electronic health record (EHR) field
that can be rendered as a sentence.
This makes the approach readily transferable to other structured clinical
datasets without requiring a fine-tuned clinical language model.

Several limitations should be considered.
First, this was a single-institution retrospective cohort, which may limit
generalisability and introduce selection bias.
Second, the non-significant AUC improvement reflects the modest sample size
($n = 247$) rather than a definitive negative result, and larger
multi-institutional cohorts are needed to confirm the observed trends.
Third, the sentence encoder was applied in a zero-shot manner without
fine-tuning on NSCLC-specific clinical text, which may limit its sensitivity
to domain-specific terminology.
Finally, PD-L1 TPS was scored using the Sauter method in this cohort, which
may differ from the SP142 or 22C3 assays used at other institutions, and the
generalisability of PD-L1-derived features should be interpreted with this
in mind.
